\section{Discussion and Limitations}
\label{sec:discussion}

\textbf{Optimizer interfaces as semantic interfaces.}  Public optimizer claims should name action, guard, layer, placement, decision time, and evidence surface.  These fields do not expose private cost formulas; they keep adapters, remote sources, and service stages from collapsing into one phenomenon, a boundary visible in federation~\cite{carey1995garlic} and polystores~\cite{stonebraker2015bigdawg}.

\textbf{What changes in an evaluation.}  Before ranking latency, cardinality error, or plan quality, a study can test whether compared systems share the same public contract relation.  If a pushdown table equates connector delegation with runtime filtering, the study can retain separating fields, narrow the claim, or report mixed phenomena.

\textbf{Using certificates with existing benchmarks.}  JOB, TPC-H, and TPC-DS standardize inputs and metrics, not contract equality.  A study can still execute those workloads; OptSem-C states which optimizer contracts the runs may compare and which layer, placement, or evidence surface must be separated.

\textbf{Protocol, scope, and refresh.}  Every optimizer comparison is a projection: choose a vocabulary, extract public contracts, validate the denominator, compute the kernel, and publish retained fields.  Counts are lower bounds over public evidence; later evidence becomes a source refresh.

\section{Conclusion}
OptSem-C turns coarse optimizer labels into public-evidence certificates that expose the contracts being equated, the fields needed to repair the denominator, and the counter-witnesses that break the abstraction.  Across seven analytical engines, layer, placement, decision time, observability, and action variant must be validated before performance claims are ranked.
